\begin{ex}
Có 6 quả táo chia đều cho 2 bạn. Hỏi mỗi bạn được bao nhiêu quả táo?
\shortans{3}
\loigiai{
Mỗi bạn được: $6 : 2 = 3$ quả táo.
}
\end{ex}

\begin{ex}
Có 8 viên bi chia đều cho 2 bạn. Hỏi mỗi bạn được bao nhiêu viên bi?
\shortans{4}
\loigiai{
Mỗi bạn được: $8 : 2 = 4$ viên bi.
}
\end{ex}

\begin{ex}
Có 10 cái kẹo chia đều vào 2 đĩa. Hỏi mỗi đĩa có bao nhiêu cái kẹo?
\shortans{5}
\loigiai{
Mỗi đĩa có: $10 : 2 = 5$ cái kẹo.
}
\end{ex}

\begin{ex}
Có 15 bông hoa cắm đều vào 5 lọ hoa. Hỏi mỗi lọ có bao nhiêu bông hoa?
\shortans{3}
\loigiai{
Mỗi lọ có: $15 : 5 = 3$ bông hoa.
}
\end{ex}

\begin{ex}
Có 20 chiếc bút chì chia đều cho 5 bạn. Hỏi mỗi bạn được bao nhiêu chiếc bút chì?
\shortans{4}
\loigiai{
Mỗi bạn được: $20 : 5 = 4$ chiếc bút.
}
\end{ex}

\begin{ex}
Có 12 cái bánh chia đều cho 2 bạn. Mỗi bạn được bao nhiêu cái bánh?
\shortans{6}
\loigiai{
Mỗi bạn được: $12 : 2 = 6$ cái bánh.
}
\end{ex}

\begin{ex}
Có 14 quyển vở chia đều cho 2 tổ. Mỗi tổ được bao nhiêu quyển vở?
\shortans{7}
\loigiai{
Mỗi tổ được: $14 : 2 = 7$ quyển vở.
}
\end{ex}

\begin{ex}
Có 25 lít dầu chia đều vào 5 can. Mỗi can có bao nhiêu lít dầu?
\shortans{5}
\loigiai{
Mỗi can có: $25 : 5 = 5$ lít dầu.
}
\end{ex}

\begin{ex}
Có 30 kg gạo chia đều vào 5 túi. Mỗi túi nặng bao nhiêu ki-lô-gam?
\shortans{6}
\loigiai{
Mỗi túi nặng: $30 : 5 = 6$ kg.
}
\end{ex}

\begin{ex}
Từ phép nhân $2 \times 6 = 12$, viết phép chia $12 : 2$. Kết quả bằng bao nhiêu?
\shortans{6}
\loigiai{
Ta có $12 : 2 = 6$.
}
\end{ex}

\begin{ex}
Từ phép nhân $5 \times 4 = 20$, viết phép chia $20 : 5$. Kết quả bằng bao nhiêu?
\shortans{4}
\loigiai{
Ta có $20 : 5 = 4$.
}
\end{ex}

\begin{ex}
Từ phép nhân $2 \times 8 = 16$, viết phép chia $16 : 8$. Kết quả bằng bao nhiêu?
\shortans{2}
\loigiai{
Ta có $16 : 8 = 2$.
}
\end{ex}

\begin{ex}
Từ phép nhân $5 \times 7 = 35$, viết phép chia $35 : 7$. Kết quả bằng bao nhiêu?
\shortans{5}
\loigiai{
Ta có $35 : 7 = 5$.
}
\end{ex}

\begin{ex}
Có 18 học sinh xếp thành các hàng, mỗi hàng có 2 bạn. Hỏi xếp được bao nhiêu hàng?
\shortans{9}
\loigiai{
Số hàng: $18 : 2 = 9$ hàng.
}
\end{ex}

\begin{ex}
Có 40 bạn học sinh xếp thành các nhóm, mỗi nhóm có 5 bạn. Hỏi có bao nhiêu nhóm?
\shortans{8}
\loigiai{
Số nhóm: $40 : 5 = 8$ nhóm.
}
\end{ex}

\begin{ex}
Có 16 chiếc ghế xếp đều vào 2 dãy. Mỗi dãy có bao nhiêu chiếc ghế?
\shortans{8}
\loigiai{
Mỗi dãy có: $16 : 2 = 8$ chiếc ghế.
}
\end{ex}

\begin{ex}
Có 35 củ cà rốt chia đều cho 5 chú thỏ. Mỗi chú thỏ được bao nhiêu củ cà rốt?
\shortans{7}
\loigiai{
Mỗi chú thỏ được: $35 : 5 = 7$ củ cà rốt.
}
\end{ex}

\begin{ex}
Một sợi dây dài 10 dm cắt thành 2 đoạn bằng nhau. Mỗi đoạn dài bao nhiêu đề-xi-mét?
\shortans{5}
\loigiai{
Mỗi đoạn dài: $10 : 2 = 5$ dm.
}
\end{ex}

\begin{ex}
Một thanh gỗ dài 45 cm cắt thành 5 đoạn bằng nhau. Mỗi đoạn dài bao nhiêu xăng-ti-mét?
\shortans{9}
\loigiai{
Mỗi đoạn dài: $45 : 5 = 9$ cm.
}
\end{ex}

\begin{ex}
Có 18 con cá thả đều vào 2 bể kính. Mỗi bể kính có bao nhiêu con cá?
\shortans{9}
\loigiai{
Mỗi bể có: $18 : 2 = 9$ con cá.
}
\end{ex}

\begin{ex}
Có 50 lít nước mắm đóng đều vào các can 5 lít. Hỏi đóng được bao nhiêu can?
\shortans{10}
\loigiai{
Đóng được: $50 : 5 = 10$ can.
}
\end{ex}

\begin{ex}
Có 20 cái kẹo chia cho mỗi bạn 2 cái. Hỏi chia được cho bao nhiêu bạn?
\shortans{10}
\loigiai{
Chia được cho: $20 : 2 = 10$ bạn.
}
\end{ex}

\begin{ex}
Điền số thích hợp: $14 : 2 = \dots$
\shortans{7}
\loigiai{
Ta có $14 : 2 = 7$.
}
\end{ex}

\begin{ex}
Điền số thích hợp: $45 : 5 = \dots$
\shortans{9}
\loigiai{
Ta có $45 : 5 = 9$.
}
\end{ex}

\begin{ex}
Tính: $20 : 2 - 3$.
\shortans{7}
\loigiai{
Ta có $20 : 2 = 10$, rồi $10 - 3 = 7$.
}
\end{ex}
